% ============================================================================
% PREAMBLE COMÚN PARA DOCUMENTOS BOOKMATE
% Curso: CC341 - Ingeniería de Software
% Universidad Nacional de Ingeniería (UNI)
% Grupo: 6
% ============================================================================

% ---------- Clase de documento ----------
\documentclass[12pt,a4paper,oneside]{article}

% ---------- Codificación y lenguaje ----------
\usepackage[utf8]{inputenc}
\usepackage[spanish,es-tabla]{babel}
\usepackage[T1]{fontenc}

% ---------- Geometría de página ----------
\usepackage[
    top=2.5cm,
    bottom=2.5cm,
    left=3cm,
    right=2.5cm,
    headheight=15pt
]{geometry}

% ---------- Fuentes ----------
\usepackage{lmodern}
\usepackage{microtype}

% ---------- Colores ----------
\usepackage[dvipsnames,table]{xcolor}
\definecolor{uniblue}{RGB}{0,51,102}
\definecolor{codegreen}{RGB}{0,128,0}
\definecolor{codegray}{RGB}{128,128,128}
\definecolor{codepurple}{RGB}{170,0,170}
\definecolor{backcolour}{RGB}{245,245,245}
\definecolor{linkcolor}{RGB}{0,0,139}

% ---------- Gráficos e imágenes ----------
\usepackage{graphicx}
\usepackage{float}
\usepackage{wrapfig}
\graphicspath{{./images/}{../commons/images/}}

% ---------- Tablas ----------
\usepackage{array}
\usepackage{longtable}
\usepackage{booktabs}
\usepackage{multirow}
\usepackage{multicol}
\usepackage{tabularx}
\usepackage{colortbl}

% ---------- Listas ----------
\usepackage{enumitem}
\setlist{noitemsep,topsep=5pt}

% ---------- Enlaces e hipervínculos ----------
\usepackage{hyperref}
\hypersetup{
    colorlinks=true,
    linkcolor=uniblue,
    filecolor=magenta,
    urlcolor=linkcolor,
    citecolor=uniblue,
    pdftitle={BookMate - CC341},
    pdfauthor={Grupo 6.2},
    pdfsubject={Ingeniería de Software},
    pdfkeywords={BookMate, Recomendación de Libros, IA},
    bookmarksnumbered=true,
    bookmarksopen=true,
    pdfstartview=FitH
}
\usepackage{bookmark}

% ---------- Bibliografía con BibLaTeX ----------
\usepackage[
    backend=biber,
    style=ieee,
    sorting=none,
    maxbibnames=99,
    doi=true,
    url=true,
    isbn=true
]{biblatex}
\addbibresource{../commons/referencias.bib}

% ---------- Código fuente ----------
\usepackage{listings}
\lstdefinestyle{mystyle}{
    backgroundcolor=\color{backcolour},
    commentstyle=\color{codegreen},
    keywordstyle=\color{blue}\bfseries,
    numberstyle=\tiny\color{codegray},
    stringstyle=\color{codepurple},
    basicstyle=\ttfamily\small,
    breakatwhitespace=false,
    breaklines=true,
    captionpos=b,
    keepspaces=true,
    numbers=left,
    numbersep=8pt,
    showspaces=false,
    showstringspaces=false,
    showtabs=false,
    tabsize=2,
    frame=single,
    rulecolor=\color{black!30}
}
\lstset{style=mystyle}

% Configurar lenguajes
\lstdefinelanguage{json}{
    basicstyle=\ttfamily\small,
    numbers=left,
    numberstyle=\scriptsize,
    stepnumber=1,
    numbersep=8pt,
    showstringspaces=false,
    breaklines=true,
    frame=single,
    string=[s]{"}{"},
    stringstyle=\color{blue},
    comment=[l]{:},
    commentstyle=\color{black},
}

% ---------- Diagramas con TikZ ----------
\usepackage{tikz}
\usetikzlibrary{shapes,arrows,positioning,calc,fit,backgrounds,shadows}

% ---------- Cajas y entornos ----------
\usepackage[most]{tcolorbox}
\tcbuselibrary{listings,breakable}

% Definir cajas personalizadas
\newtcolorbox{infobox}[1][]{
    colback=blue!5!white,
    colframe=uniblue,
    fonttitle=\bfseries,
    title=Información,
    #1
}

\newtcolorbox{warningbox}[1][]{
    colback=yellow!10!white,
    colframe=orange,
    fonttitle=\bfseries,
    title=Advertencia,
    #1
}

\newtcolorbox{notebox}[1][]{
    colback=green!5!white,
    colframe=green!75!black,
    fonttitle=\bfseries,
    title=Nota,
    #1
}

% ---------- Matemáticas ----------
\usepackage{amsmath,amssymb,amsfonts}
\usepackage{mathtools}

% ---------- Headers y footers ----------
\usepackage{fancyhdr}
\pagestyle{fancy}
\fancyhf{}
\fancyhead[L]{\small\textit{\docTitle}}
\fancyhead[R]{\small CC341 - Grupo 6.2}
\fancyfoot[C]{\thepage}
\renewcommand{\headrulewidth}{0.4pt}
\renewcommand{\footrulewidth}{0.4pt}

% Para primera página sin header
\fancypagestyle{plain}{
    \fancyhf{}
    \renewcommand{\headrulewidth}{0pt}
    \fancyfoot[C]{\thepage}
}

% ---------- Títulos de secciones ----------
\usepackage{titlesec}
\titleformat{\section}
    {\color{uniblue}\normalfont\Large\bfseries}
    {\thesection}{1em}{}
\titleformat{\subsection}
    {\color{uniblue}\normalfont\large\bfseries}
    {\thesubsection}{1em}{}
\titleformat{\subsubsection}
    {\color{uniblue}\normalfont\normalsize\bfseries}
    {\thesubsubsection}{1em}{}

% ---------- Espaciado ----------
\usepackage{setspace}
\setstretch{1.15}
\setlength{\parindent}{0pt}
\setlength{\parskip}{6pt}

% ---------- Caption ----------
\usepackage[font=small,labelfont=bf]{caption}
\usepackage{subcaption}

% ---------- Acrónimos y glosario ----------
\usepackage[acronym,toc]{glossaries}

% ---------- Unidades SI ----------
\usepackage{siunitx}

% ---------- Referencias cruzadas ----------
\usepackage{cleveref}
\crefname{figure}{Figura}{Figuras}
\crefname{table}{Tabla}{Tablas}
\crefname{section}{Sección}{Secciones}

% ---------- PlantUML integration ----------
\usepackage{pdfpages}

% ---------- PDF features ----------
\usepackage{pdfpages}

% ---------- Metadatos adicionales ----------
\usepackage{datetime}
\newdateformat{mydate}{\THEDAY\ de \monthname[\THEMONTH] de \THEYEAR}

% ============================================================================
% COMANDOS PERSONALIZADOS
% ============================================================================

% Comando para portada estándar UNI
\newcommand{\makeUNIcover}{
    \begin{titlepage}
        \centering
        
        % Logo UNI (si está disponible)
        % \includegraphics[width=4cm]{logo-uni.png}\\[1cm]
        
        {\scshape\Large Universidad Nacional de Ingeniería\par}
        {\scshape\large Facultad de Ingeniería Industrial y de Sistemas\par}
        \vspace{0.5cm}
        {\scshape\normalsize Escuela Profesional de Ingeniería de Sistemas\par}
        \vspace{1.5cm}
        
        {\huge\bfseries\color{uniblue} \docTitle \par}
        \vspace{0.5cm}
        {\Large\itshape \docSubtitle \par}
        \vspace{2cm}
        
        {\large \textbf{Curso:} CC341 - Ingeniería de Software\par}
        \vspace{0.3cm}
        {\large \textbf{Proyecto:} BookMate - Sistema de Recomendación de Libros\par}
        \vspace{0.3cm}
        {\large \textbf{Grupo:} 6.2\par}
        \vspace{1.5cm}
        
        {\large \textbf{Integrantes:}\par}
        \vspace{0.5cm}
        \begin{tabular}{ll}
            \textbullet\ Integrante 1 & \textit{(Código)} \\
            \textbullet\ Integrante 2 & \textit{(Código)} \\
            \textbullet\ Integrante 3 & \textit{(Código)} \\
            \textbullet\ Integrante 4 & \textit{(Código)} \\
            \textbullet\ Integrante 5 & \textit{(Código)} \\
            \textbullet\ Integrante 6 & \textit{(Código)} \\
        \end{tabular}
        
        \vfill
        
        {\large \textbf{Semana:} \docWeek\par}
        {\large \textbf{Fecha:} \mydate\today\par}
        {\large \textbf{Versión:} \docVersion\par}
        
        \vspace{1cm}
        {\large Ciclo 2024-2\par}
        {\large Lima, Perú\par}
        
    \end{titlepage}
    \newpage
}

% Comando para tabla de versiones
\newcommand{\versionTable}{
    \section*{Control de Versiones}
    \begin{table}[H]
        \centering
        \begin{tabular}{|c|c|p{6cm}|p{3cm}|}
            \hline
            \textbf{Versión} & \textbf{Fecha} & \textbf{Descripción} & \textbf{Autor(es)} \\
            \hline
            0.1 & \today & Versión inicial del documento & Equipo 6.2 \\
            \hline
            & & & \\
            \hline
            & & & \\
            \hline
        \end{tabular}
        \caption{Historial de versiones}
    \end{table}
    \newpage
}

% Comando para sección de futuro
\newcommand{\futureSection}[1]{
    \textcolor{blue}{\textit{Se #1}}
}

% Comando para resaltar términos
\newcommand{\term}[1]{\textbf{\textit{#1}}}

% Comando para código inline
\newcommand{\code}[1]{\texttt{\small #1}}

% Comando para URLs
\newcommand{\weblink}[2]{\href{#1}{\textcolor{linkcolor}{#2}}}

% ============================================================================
% CONFIGURACIÓN FINAL
% ============================================================================

% Evitar widows y orphans
\widowpenalty=10000
\clubpenalty=10000

% Profundidad de tabla de contenidos
\setcounter{tocdepth}{3}
\setcounter{secnumdepth}{3}

% ============================================================================
% FIN DEL PREAMBLE
% ============================================================================
